\section{Related Works}
\textbf{Prompt Synthesis.} Research on prompt synthesis has emerged in tandem with the rapid advancement and widespread adoption of LLMs. Early studies investigated general-purpose methods. For instance, \citet{wang2023self} proposed generating new prompts with GPT-3 by instructing the model to imitate a small set of seed instructions, while \citet{xu2024wizardlm} introduced a self-evolution algorithm that increases the complexity of seed instructions through prompting ChatGPT. Subsequent work shifted toward synthesizing domain-specific prompts aimed at cultivating specialized LLM capabilities. Along this line, \citet{luo2023wizardmath} extended the self-evolution framework to mathematical reasoning, and \citet{huang2025key} as well as \citet{tang2024mathscale} proposed generating new math problems from key points or concepts. Most recently, \citet{zhao2025promptcot} incorporated a ``thinking process'' into math problem generation, further enhancing the difficulty of the synthesized problems. Beyond mathematics, prompt synthesis has also been applied to coding problem generation. For example, \citet{huang2024opencoder} and \citet{wei2023magicoder} explored generating programming problems from code snippets. More recently, with the growing interest in agentic intelligence, research on prompt synthesis has begun to evolve toward generation of task-oriented problems as a means of overcoming data scarcity in the development of LLM agents. For instance, \citet{li2025websailor} proposed synthesizing complex knowledge-seeking questions from knowledge graphs; \citet{liu2025webexplorer} developed an iterative question evolution method to progressively increase question difficulty; and \citet{yang2025swe} combined LLMs-based and rule-based techniques to create bugs in code derived from selected GitHub repositories.
A key limitation of prior work is its reliance on heuristic or single-pass generation pipelines, which often fail to capture the deeper reasoning structures needed for truly challenging problems.  
In contrast, our approach frames rationale–prompt co-generation as an EM procedure: the E-step refines rationales through reward-guided inference, while the M-step updates prompt synthesis conditioned on these rationales.  
This iterative loop provides a principled mechanism for aligning concepts, rationales, and prompts, leading to more faithful and more difficult problems in both mathematics and programming domains.


\textbf{Self-Evolution.} Progress in prompt synthesis has also paved the way for developing self-evolutionary LLMs, where prompt generation and model updating are coupled in an iterative cycle that allows models to learn from their own experience. As a promising alternative to human-curated data, such self-evolving mechanisms have attracted growing attention within the research community. Representative efforts include \citet{yuan2024self}, who leveraged the self-instruct framework \citep{wang2023self} for prompt synthesis and subsequently improved a base model via supervised fine-tuning, and \citet{zhao2025absolute}, who formalized three typical reasoning paradigms as programming tasks and iteratively combined task synthesis with reinforcement learning. PromptCoT~2.0 adopts a self-play paradigm for post-training LLMs, sharing similarities with these self-evolutionary studies. The key distinction, however, is that PromptCoT~2.0 does not rely on self-judgment, as in \citet{yuan2024self}, nor on the joint optimization of prompt synthesis and model post-training, as in \citet{zhao2025absolute}. While primarily focused on advancing prompt synthesis, PromptCoT~2.0—by virtue of the quality of its prompts and its state-of-the-art performance across diverse downstream tasks—opens avenues for future exploration along the self-evolutionary line.





